% Part: first-order-logic
% Chapter: proof-systems
% Section: axiomatic-deduction

\documentclass[../../../include/open-logic-section]{subfiles}

\begin{document}

\iftag{FOL}
      {\olfileid{fol}{prf}{axd}}
      {\olfileid{pl}{prf}{axd}}

\olsection{స్వీకృతాధారిత \usetoken{P}{derivation}}

స్వీకృతాధారిత \tetoken{వ్యుత్పత్తులు}{!!{derivation}s} అత్యంత పాతవి, అత్యంత సరళమైన తార్కిక
\tetoken{వ్యుత్పత్తి}{!!{derivation}} వ్యవస్థలు. వాటిలోని \tetoken{వ్యుత్పత్తులు}{!!{derivation}s} కేవలం
\tetoken{వాక్యాలు}{!!{sentence}s} క్రమాలు. ఒక \tetoken{వాక్యాలు}{!!{sentence}s} క్రమంలోని ప్రతి
\tetoken{వాక్యం}{!!{sentence}}~$!A$ కింది షరతుల్లో ఒకదాన్ని సంతృప్తిపరిస్తే ఆ క్రమం
సరైన \tetoken{వ్యుత్పత్తి}{!!{derivation}} అవుతుంది:
\begin{enumerate}
\item $!A$ ఒక స్వీకృతం, లేదా
\item $!A$, ఇచ్చిన \tetoken{వాక్యాలు}{!!{sentence}s} సమితి~$\Gamma$లో ఒక \tetoken{మూలకం}{!!a{element}}, లేదా
\item ఒక నిగమన నియమం $!A$ను సమర్థిస్తుంది.
\end{enumerate}
స్వీకృతం కావడానికి $!A$, స్థిరమైన కొన్ని \tetoken{వాక్యం}{!!{sentence}} నమూనాల్లో
ఒకదాని రూపంలో ఉండాలి. మొదటిస్థాయి విధేయ తర్కానికి సంతృప్తికరమైన
(నిర్దుష్టమైన, సంపూర్ణమైన) \tetoken{వ్యుత్పత్తి}{!!{derivation}} వ్యవస్థను ఇచ్చే స్వీకృత
నమూనాల సమితులు అనేకం ఉన్నాయి. కొన్ని, అవి నియంత్రించే సంయోజకాల
ప్రకారం వ్యవస్థీకరించబడ్డాయి. ఉదాహరణకు,
\[
!A \lif (!B \lif !A) \qquad !B \lif (!B \lor !C) \qquad (!B \land !C) \lif !B
\]
అనే నమూనాలు $\lif$, $\lor$, $\land$లను నియంత్రించే సాధారణ
స్వీకృతాలు. కొన్ని స్వీకృత వ్యవస్థలు స్వీకృతాల సంఖ్యను కనిష్ఠం
చేయాలని లక్ష్యంగా పెట్టుకుంటాయి. ఏ సంయోజకాలను ఆదిమాలుగా
తీసుకుంటామనే దానిని బట్టి, ఒకే స్వీకృతంతో కూడిన స్వీకృత వ్యవస్థలను
కూడా కనుగొనవచ్చు.

నిగమన నియమం అంటే, \tetoken{ఒక వ్యుత్పత్తిలో}{!!a{derivation}} ఒక \tetoken{వాక్యానికి}{!!a{sentence}} సమర్థన
ఉండటానికి సరిపోయే షరతును ఇచ్చే షరతీయ ప్రకటన. మోడస్ పోనెన్స్ ఎంతో
సాధారణమైన అలాంటి నియమం: $!A$, $!A \lif !B$లకు ఇప్పటికే సమర్థన
ఉంటే $!B$కు సమర్థన ఉందని అది చెబుతుంది. అంటే, ఏదో
\tetoken{వాక్యం}{!!{sentence}}~$!A$కు $!A$, $!A \lif !B$ రెండూ \tetoken{వ్యుత్పత్తిలో}{!!{derivation}}
$!B$కంటే ముందు వస్తే, ఆ \tetoken{ఒక వ్యుత్పత్తిలో}{!!a{derivation}} \tetoken{వాక్యం}{!!{sentence}}~$!B$ ఉన్న
వరుసకు సమర్థన ఉంటుంది.

స్వీకృతాధారిత \tetoken{వ్యుత్పత్తులు}{!!{derivation}s} ఆధారంగా $\Proves$ సంబంధాన్ని ఇలా
నిర్వచిస్తాం: చివరి సూత్రంగా \tetoken{వాక్యం}{!!{sentence}}~$!A$ ఉన్న ఒక
\tetoken{వ్యుత్పత్తి}{!!a{derivation}} ఉంటే, అప్పుడు మాత్రమే $\Gamma \Proves !A$; ఇక్కడ
పై (2) ద్వారా సమర్థన పొందిన ఆ \tetoken{వ్యుత్పత్తిలోని}{!!{derivation}} \tetoken{వాక్యాలు}{!!{sentence}s}
సమితిగా $\Gamma$ను తీసుకుంటాం. $\Gamma$ శూన్యంగా ఉండే, అంటే
\tetoken{వ్యుత్పత్తిలోని}{!!{derivation}} ప్రతి \tetoken{వాక్యం}{!!{sentence}} (1) లేదా (3) ద్వారా
సమర్థన పొందే, ఒక \tetoken{వ్యుత్పత్తి}{!!a{derivation}} $!A$కు ఉంటే $!A$ ఒక సిద్ధాంతం.
ఉదాహరణకు, $\Proves !A \lif (!B \lif (!B \lor !A))$ అని చూపే
ఒక \tetoken{వ్యుత్పత్తి}{!!a{derivation}} ఇది:
\begin{derivation}
  1. & $!B \lif (!B \lor !A)$ \\
  2. & $(!B \lif (!B \lor !A)) \lif (!A  \lif (!B \lif (!B \lor !A)))$\\
  3. & $!A  \lif (!B \lif (!B \lor !A))$
\end{derivation}
1వ వరుసలోని \tetoken{వాక్యం}{!!{sentence}}, $!A \lif (!A \lor !B)$ అనే స్వీకృత
రూపంలో ఉంది ($!A$, $!B$ల పాత్రలు తారుమారు అయ్యాయి). 2వ వరుసలోని
వాక్యం $!A \lif (!B \lif !A)$ అనే స్వీకృత రూపంలో ఉంది. కాబట్టి
ఈ రెండు వరుసలకూ సమర్థన ఉంది. 3వ వరుసకు మోడస్ పోనెన్స్ ద్వారా
సమర్థన ఉంది: దానిని $!D$గా సంక్షిప్తీకరిస్తే, 2వ వరుస
$!C \lif !D$ రూపంలో ఉంటుంది; ఇక్కడ $!C$ అంటే 1వ వరుస అయిన
$!B \lif (!B \lor !A)$.

$\Gamma \Proves \lfalse$ అయితే సమితి $\Gamma$ వైరుధ్యభరితమైనది.
సంపూర్ణమైన స్వీకృత వ్యవస్థ ఏ $!A$కైనా $\lfalse \lif !A$ అని కూడా
నిరూపిస్తుంది. అందువల్ల $\Gamma$ వైరుధ్యభరితమైతే, ఏ $!A$కైనా
$\Gamma \Proves !A$.

తర్కానికి స్వీకృతాధారిత \tetoken{వ్యుత్పత్తులు}{!!{derivation}s} వ్యవస్థలను మొదట గాట్లోబ్
ఫ్రేగె తన 1879 నాటి \emph{Begriffsschrift}లో ఇచ్చాడు; ఈ కారణంగానే
దానిని ఆధునిక తర్కశాస్త్రపు మొదటి రచనగా తరచుగా పరిగణిస్తారు.
ఆల్ఫ్రెడ్ నార్త్ వైట్‌హెడ్, బెర్ట్రాండ్ రసెల్‌ల
\emph{Principia Mathematica}లోనూ, 1920లలో డేవిడ్ హిల్బర్ట్, అతని
విద్యార్థుల చేతిలోనూ అవి పరిపూర్ణ రూపం పొందాయి. అందువల్ల వాటిని
“ఫ్రేగె వ్యవస్థలు” లేదా “హిల్బర్ట్ వ్యవస్థలు” అని తరచుగా అంటారు.
ఒక తర్కవ్యవస్థకు స్వీకృతాధారిత వ్యవస్థను కనుగొనడం తరచుగా సులభం
కనుక అవి ఎంతో అనువైనవి. \tetoken{వ్యుత్పత్తులు}{!!{derivation}s} నిర్మాణం చాలా సరళంగా,
నిగమన నియమాలు ఒకటో రెండో మాత్రమే ఉంటాయి కనుక వాటి
\emph{గురించి} విషయాలను నిరూపించడం కూడా సాపేక్షంగా సులభం. అయితే,
ఆచరణలో వాటిని ఉపయోగించడం చాలా కష్టం; అంటే నిరూపణలను కనుగొని
రాయడం కష్టం.

\end{document}
